\documentclass[12pt]{article}

\usepackage[a4paper, margin=1in]{geometry}
\usepackage{graphicx}
\usepackage{booktabs}
\usepackage{amsmath}
\usepackage{hyperref}
\usepackage{float}

\title{\textbf{Banking Customer Churn Prediction}\\
	\large A comparative analysis of various supervised ML models \\to find the most suitable one resilient to severe class imbalance} 
\author{Harshank Nimonkar}
\date{\vspace{-5ex}}

\begin{document}
	
	\maketitle
	
	\begin{abstract}
		\noindent This report presents a comparative analysis of supervised machine learning models for customer churn prediction. The dataset is highly imbalanced (approximately 95\% non-churn vs 5\% churn), making accuracy an unreliable performance indicator. Accordingly, models are evaluated using imbalance-sensitive metrics such as Precision, Recall, F1-score, ROC-AUC, and Precision–Recall curves. Empirical results suggest that ensemble methods provide improved minority-class detection and stronger ranking performance relative to linear approaches, offering more effective identification of high-risk customers.
	\end{abstract}
	
	\section{Problem Statement}
	
	The objective of this study is to predict whether a customer will churn based on historical behavioral and transactional features. 
	Given the highly imbalanced nature of the dataset, the primary goal is to maximize minority-class detection performance while maintaining reasonable precision.
	
	\section{Dataset Characteristics}
	
	\begin{itemize}
		\item \href{https://huggingface.co/datasets/electricsheepafrica/nigerian-banking-customer-360}{Nigerian Banking Customer 360}	
		\item 22 original features and 2 targets (30d churn and 90d churn).
		\item For this analysis, we choose the 30 day churn target variable.
		\item 2 polytomous features with 3 and 37 categories; applied OneHot encoding.
		\item Total samples: 2M
		\item Majority class (Non-Churn): $\sim$95\%
		\item Minority class (Churn): $\sim$5\%
		\item Expanded feature space to 60 dimensions
	\end{itemize}
	
	A stratified train-test split (70\%-30\%) was used to preserve class distribution.
	
	\section{Evaluation Strategy}
	
	Due to class imbalance, the following metrics were prioritized:
	
	\begin{itemize}
		\item Minority Precision
		\item Minority Recall
		\item Minority F1-score
		\item ROC Curve and AUC
		\item Precision-Recall Curve and Average Precision (AP)
	\end{itemize}
	
	Accuracy is reported but not used for model selection.
	
	\section{Models Evaluated}
	
	The following models were compared:
	
	\begin{enumerate}
		\item Logistic Regression (Baseline) (LR)
		\item Bernoulli Naive Bayes (BNB)
		\item SGD Classifier (solver='hinge') (SGD)
		\item Decision Tree (DT)
		\item Random Forest (RF)
		\item AdaBoost (AB)
		\item Gradient Boosting (GB)
	\end{enumerate}
	
	All models were trained under identical preprocessing and evaluation conditions.
	
	\begin{table}[H]
		\centering
		\begin{tabular}{lccccccc}
			\toprule
			Model & Accuracy & '1' Precision & '0' Precision & '1' Recall & '0' Recall & '1' F1 & '0' F1 \\
			\midrule
			LR & 0.9466 & 0.4101 & 0.9569 & 0.1556 & 0.9882 & 0.2256 & 0.9723\\
			BNB & 0.9502 & 0.6751 & 0.9503 & 0.0062 & 0.9998 & 0.0124 & 0.9744\\
			SGD  & 0.9481 & 0.3593 & 0.9521 & 0.0492 & 0.9954 & 0.0866 & 0.9733\\ 
			DT  & 0.9904 & 0.8950 & 0.9955 & 0.9153 & 0.9943 & 0.9051 & 0.9949\\ 
			RF  & 0.9904 & 0.8951 & 0.9955 & 0.9154 & 0.9944 & 0.9051 & 0.9949\\
			AB  & 0.9850 & 0.8939 & 0.9893 & 0.7950 & 0.9950 & 0.8416 & 0.9921\\ 
			GB  & 0.9904 & 0.8950 & 0.9955 & 0.9154 & 0.9943 & 0.9051 & 0.9949\\
			\bottomrule
		\end{tabular}
		\caption{Performance of different models evaluated across different metrics such as Precision, Recall and F1 scores for the majority and the minority classes. Class '1': Minority Class (Churn), Class '0': Majority Class (No Churn).}
	\end{table}
	
	\section{Baseline Performance}
	
	Logistic Regression was used as a baseline linear model.
	
	\begin{itemize}
		\item Accuracy $\approx$ 94.66\%
		\item Minority Recall relatively low
		\item Minority F1 significantly lower than overall accuracy
	\end{itemize}
	
	This demonstrates that accuracy is misleading in imbalanced classification problems.
	
	\section{Tree-Based Models}
	
	The Decision Tree model showed substantial performance improvement:
	
	\begin{itemize}
		\item Significant increase in Minority Recall
		\item Strong improvement in Minority F1-score
		\item Overall accuracy $\approx$ 99\%
	\end{itemize}
	
	This indicates nonlinear interactions among features.
	
	\section{Ensemble Methods}
	
	\subsection{Random Forest}
	
	Random Forest reduced variance and stabilized tree performance:
	
	\begin{itemize}
		\item High minority detection capability
		\item Strong macro and weighted F1 scores
	\end{itemize}
	
	\subsection{AdaBoost}
	
	AdaBoost improved Churn recall relative to baseline but was significantly less consistent than Random Forest and Gradient Boosting. 
	
	\subsection{Gradient Boosting}
	
	Gradient Boosting demonstrated the strongest and most consistent performance:
	
	\begin{itemize}
		\item Highest Minority F1-score
		\item Highest Average Precision
		\item Strong ROC and PR curve dominance
	\end{itemize}
	
	\section{Minority Class Analysis: Precision vs Recall}
	
	\begin{figure}[H]
		\centering
		\includegraphics[width=\linewidth]{Minority_Precision_vs_Recall.png}
		\caption{Minority Precision vs Minority Recall. Each star represents an individual model.}
		\label{fig:pvr1}
	\end{figure}
		
	\begin{figure}[H]
	\centering
	\includegraphics[width=\linewidth]{Minority_Precision_vs_Recall2.png}
	\caption{Minority Precision vs Minority Recall for top 3 performing models.}
	\label{fig:pvr2}
	\end{figure}
	
		The precision and recall for the Churn class for different models shows a clear separation between the linear and non-linear models. The SGD (with 'hinge' loss) is basically a Support Vector Classifier. The Bernoulli Naive Bayes Classifier, though a better guessor than LR and SGD, still fails to capture the non-linearities in the features. Tree-based and ensemble models succeed in this regard as evident from Fig. \ref{fig:pvr1}. We get a better look at the top right cluster in Fig. \ref{fig:pvr2}.\ 
	\newline Looking at Fig. \ref{fig:pvr2}, the Precision-Recall tradeoff is apparent from the RF and GB positions. Depending on the business cost trade-offs, if retention outreach is inexpensive relative to churn loss, the model with higher recall will be preferred, i.e., Gradient Boosting. If outreach is expensive or limited in capacity, higher precision will be preferred, i.e., Random Forest.
	
		
	\section{ROC Curve Analysis}
	
	Figure~\ref{fig:roc} shows the ROC curve comparison across models. While all models exceed the random baseline, the ROC-AUC may overstate performance in this highly imbalanced setting (~5\% minority class), since the false positive rate is normalized by the large majority class. Consequently, even a considerable number of false positives can correspond to a relatively small false positive rate, making the curves appear optimistic.
	
	\begin{figure}[H]
		\centering
		\includegraphics[width=\linewidth]{Figure_1.png}
		\caption{ROC Curve Comparison Across Models}
		\label{fig:roc}
	\end{figure}
	
	\section{Precision-Recall Curve Analysis}
	
	Figure~\ref{fig:pr} presents the Precision-Recall curve comparison. With a churn prevalence of approximately 5\%, the expected precision under random prediction is $\approx 0.05$, establishing a natural performance baseline. Ensemble models significantly exceed this baseline and outperform linear models in ranking minority-class instances, as reflected in improved ROC-AUC and precision-recall characteristics. This suggests superior ability to capture nonlinear interactions and better separate churners from non-churners in probability space, resulting in higher lift within top-ranked segments.
	
	\begin{figure}[H]
		\centering
		\includegraphics[width=\linewidth]{Figure_2.png}
		\caption{Precision-Recall Curve Comparison Across Models}
		\label{fig:pr}
	\end{figure}
		
	\section{Model Selection}
	
	Given the business objective of churn detection, minority recall and F1-score were prioritized. From the current analysis, Gradient Boosting is recommended closely followed by Random Forest due to:
	
	\begin{itemize}
		\item Superior minority detection
		\item Strong precision-recall tradeoff
		\item Robust overall performance
	\end{itemize}
	
	Further hyperparameter tuning may improve the precision–recall balance and support a more data-driven model selection decision. 
	
	\section{Key Takeaways}
	
	\begin{enumerate}
		\item Accuracy is insufficient in imbalanced classification.
		\item Linear models underperform in nonlinear feature spaces.
		\item Tree-based ensembles significantly improve minority detection.
		\item Precision-Recall analysis is critical for rare-event prediction.
	\end{enumerate}
	
	\section{Future Work}
	
	Future improvements may include:
	
	\begin{itemize}
		\item Hyperparameter optimization
		\item Threshold tuning for business cost optimization
		\item Feature importance analysis and interpretation
		\item Cost-sensitive learning approaches
	\end{itemize}
	
\end{document}